\subsection{Análise Comparativa Detalhada: CORA-Eval vs Estado da Arte}

\subsubsection{MATH (Hendrycks et al., 2021)}

O benchmark MATH\footnote{Hendrycks, D. et al. \textit{Measuring Mathematical
Problem Solving With the MATH Dataset}. NeurIPS, 2021. DOI:
10.48550/arXiv.2103.03874. Dataset de 12.500 problemas matemáticos de
competição (AMC, AIME) com 7 áreas e 5 níveis, disponível em
\url{https://github.com/hendrycks/math}.} é o benchmark mais próximo do
CORA-Eval em espírito avaliativo. Ambos priorizam respostas em formato
livre (não múltipla escolha) e organizam tarefas por níveis de dificuldade.
Contudo, três diferenças fundamentais limitam a capacidade do MATH de capturar
maturidade científica integrada:

\textbf{Primeiro, monotematicidade.} O MATH avalia exclusivamente matemática,
enquanto a prática científica real exige integração de múltiplos domínios. Um
químico computacional não apenas resolve equações --- modela sistemas
moleculares, interpreta espectros, otimiza geometrias. O CORA-Eval captura
esta multidimensionalidade através de 10 dimensões.

\textbf{Segundo, ausência de verificação simbólica.} A correção no MATH é
determinística: compara-se a resposta final do modelo com o gabarito. Não há
verificação de que o raciocínio intermediário é correto --- uma resposta
correta pode ser obtida por acaso, por memorização ou por um caminho lógico
inválido. O CORA-Eval exige que múltiplos verificadores independentes
(V1--V7) aprovem cada resposta.

\textbf{Terceiro, natureza estática.} O MATH é um snapshot: avalia o desempenho
em um instante, sem rastrear evolução. O CORA-Eval mantém 8 snapshots
temporais com métricas de crescimento (tensor de evolução) e projeções.

\subsubsection{MMLU (Hendrycks et al., 2020)}

O MMLU\footnote{Hendrycks, D. et al. \textit{Measuring Massive Multitask
Language Understanding}. ICLR, 2021. DOI: 10.48550/arXiv.2009.03300.
15.908 questões de múltipla escolha em 57 áreas.} cobre 57 áreas de
conhecimento --- a maior amplitude entre benchmarks estabelecidos --- mas
avalia exclusivamente conhecimento declarativo através de múltipla escolha.
Esta limitação é crítica para a avaliação de raciocínio científico: um
sistema pode selecionar a alternativa correta sem compreender o processo
que a produziu.

Evidência empírica desta limitação: quando as 150 tarefas do CORA-Eval foram
convertidas para formato de múltipla escolha (adicionando 3 distratores
plausíveis), o desempenho de modelos locais (Ollama) aumentou em média 23\%
em relação ao formato de resposta livre. Esta discrepância ($\Delta = 23\%$)
representa a diferença entre ``reconhecer a resposta correta'' e ``produzir
a resposta correta'' --- precisamente a distinção que o CORA-Eval foi projetado
para capturar.

\subsubsection{SciBench (Wang et al., 2023)}

O SciBench\footnote{Wang, X. et al. \textit{SciBench: Evaluating College-Level
Scientific Problem-Solving Abilities of Large Language Models}. arXiv:2307.10635,
2023. DOI: 10.48550/arXiv.2307.10635.} é o benchmark mais próximo do CORA-Eval
em escopo científico, com 695 problemas de livros-texto universitários em
física, química e matemática. Sua principal contribuição é a categorização
de erros em 5 tipos (conceitual, computacional, unidade, interpretação,
formulação).

O CORA-Eval estende esta categorização através dos 7 verificadores V1--V7:
erros de unidade são detectados por V1 (dimensional), erros computacionais
por V5 (numérico), erros conceituais por V3 (contraexemplos), erros de
formulação por V2 (algébrico), e erros de interpretação por V4 (estatístico).
Esta correspondência não é acidental --- valida a taxonomia de erros do
SciBench e demonstra que os verificadores Cora cobrem o espectro completo
de modos de falha em raciocínio científico.

\subsubsection{BIG-bench (Srivastava et al., 2022)}

O BIG-bench\footnote{Srivastava, A. et al. \textit{Beyond the Imitation Game:
Quantifying and Extrapolating the Capabilities of Language Models}.
arXiv:2206.04615, 2022. DOI: 10.48550/arXiv.2206.04615. Consórcio com 450+
autores, 204 tarefas.} representa o extremo oposto do CORA-Eval em filosofia
de design: máxima amplitude (204 tarefas) com mínima profundidade (cada tarefa
avaliada isoladamente, sem conexão entre domínios). Esta abordagem é valiosa
para exploração inicial de capacidades, mas não permite avaliação de
maturidade integrada.

O CORA-Eval adota a filosofia inversa: amplitude moderada (150 tarefas) com
máxima profundidade (cada tarefa verificada por 3--7 métodos independentes,
rastreável a testes TDD, validada externamente). Esta escolha reflete uma
convicção metodológica: \textbf{na avaliação de raciocínio científico, a
profundidade da verificação é mais importante que a amplitude da cobertura}.

\subsection{Teoria da Decisão aplicada à Seleção de Verificadores}

A seleção de quais verificadores Cora (V1--V7) aplicar a cada tarefa é
formalizada como um problema de decisão sequencial sob incerteza. O motor
Q-Score UCB1\footnote{Auer, P., Cesa-Bianchi, N., Fischer, P. \textit{Finite-time
Analysis of the Multi-armed Bandit Problem}. Machine Learning, 47:235--256,
2002. DOI: 10.1023/A:1013689704352.} implementa o equilíbrio exploração
\textit{vs} exploração (exploration-exploitation):

\[
Q_i = \bar{v}_i + \sqrt{\frac{2\ln N}{n_i}}
\]

onde $\bar{v}_i$ é a recompensa média histórica do verificador $i$ (taxa de
aprovação), $N$ é o número total de verificações realizadas, e $n_i$ é o
número de vezes que o verificador $i$ foi aplicado. O termo $\sqrt{2\ln N /
n_i}$ é o bônus de exploração: verificadores pouco utilizados ($n_i$ pequeno)
recebem um bônus maior, incentivando sua aplicação para reduzir a incerteza
sobre sua eficácia.

No estado atual do ecossistema, V5 (Numérico) possui $n_5 = 60+$ aplicações
e $\bar{v}_5 = 0,95$ (alta confiança), enquanto V7 (Código) possui $n_7 = 8$
e $\bar{v}_7 = 0,90$ (confiança moderada). O algoritmo UCB1, portanto,
priorizará a aplicação de V7 nas próximas avaliações para reduzir a incerteza
--- uma prescrição concreta e quantitativa para a evolução do ecossistema.

\subsection{Análise de Robustez: Perturbação dos Scores}

Para avaliar a robustez do CORA-Score a erros de medição nos scores
individuais, aplicamos ruído gaussiano $\mathcal{N}(0, \sigma^2)$ a cada
$S_d$ e recalculamos o CORA-Score para $\sigma \in \{0,05; 0,10; 0,20\}$,
com 1.000 replicações bootstrap.

Resultados: para $\sigma = 0,10$ (erro de $\pm 0,10$ em cada score, que
corresponde a aproximadamente 1 tarefa de imprecisão), a distribuição do
CORA-Score tem média 3,04 e desvio padrão 0,042. O intervalo de confiança
de 95\% é [2,96; 3,12], que não se sobrepõe ao limiar de M3 (2,50) nem ao
de M5 (4,00). Isto confirma que a classificação ``Pesquisa'' (M4) é robusta
a erros de medição moderados.

Para $\sigma = 0,20$ (erro de $\pm 0,20$, correspondente a $\sim$2 tarefas),
o IC 95\% expande para [2,88; 3,19], ainda mantendo-se acima do limiar de
M3. A classificação M4 é, portanto, robusta mesmo sob incerteza considerável
nos scores individuais.

\subsection{Distribuição de Tarefas por Dimensão e Nível: Análise de Cobertura}

\begin{table}[H]\centering\footnotesize
\caption{Distribuição de tarefas e cobertura por dimensão e nível}
\label{tab:coverage}
\begin{tabular}{p{0.05\textwidth}p{0.18\textwidth}p{0.06\textwidth}p{0.06\textwidth}p{0.06\textwidth}p{0.06\textwidth}p{0.10\textwidth}p{0.10\textwidth}}
\toprule
\textbf{D} & \textbf{Dimensão} & \textbf{N1} & \textbf{N2} & \textbf{N3} & \textbf{N4} & \textbf{Total} & \textbf{Cobertura} \\
\midrule
D1 & Matemática & 4/4 & 5/5 & 4/5 & 4/5 & 17/19 & 89,5\% \\
D2 & Física & 3/3 & 4/4 & 4/4 & 2/4 & 13/15 & 86,7\% \\
D3 & Estatística & 3/3 & 5/5 & 3/5 & 2/5 & 13/18 & 72,2\% \\
D4 & Química & 3/3 & 4/4 & 1/4 & 0/3 & 8/14 & 57,1\% \\
D5 & Biologia & 3/3 & 4/4 & 2/4 & 0/3 & 9/14 & 64,3\% \\
D6 & Geociências & 3/3 & 3/3 & 2/3 & 0/3 & 8/12 & 66,7\% \\
D7 & Código & 3/3 & 3/4 & 4/5 & 1/5 & 11/17 & 64,7\% \\
D8 & Literatura & 3/3 & 4/4 & 1/4 & 0/4 & 8/15 & 53,3\% \\
D9 & Metodologia & 2/3 & 3/4 & 3/4 & 0/4 & 8/15 & 53,3\% \\
D10 & Síntese & 2/2 & 3/3 & 3/3 & 2/3 & 10/11 & 90,9\% \\
\midrule
\multicolumn{5}{r}{\textbf{Total}} & \textbf{105/150} & \textbf{70,0\%} \\
\bottomrule
\end{tabular}
\end{table}

A cobertura geral de 70,0\% (105/150 tarefas aprovadas) é assimétrica: D1 e
D10 excedem 89\%, enquanto D8 e D9 estão abaixo de 54\%. Esta assimetria é
consistente com a arquitetura do ecossistema, otimizada para raciocínio formal
(D1, D10) e menos desenvolvida para processamento de texto em larga escala
(D8) e metodologia experimental (D9). A geodésica de aprendizado (Seção 9)
prescreve D8 e D9 como prioridades de desenvolvimento.

\subsection{Contribuições Específicas para a Ciência da Computação}

Esta dissertação oferece cinco contribuições específicas para a ciência da
computação:

\textbf{C1 --- CORA-Eval:} O primeiro benchmark evolutivo multidimensional com
verificação simbólica integrada para avaliação de maturidade científica em
ecossistemas multiagente. Diferente de benchmarks estáticos (MATH, MMLU), o
CORA-Eval rastreia a evolução temporal através de snapshots, permitindo análise
de tendências e projeções.

\textbf{C2 --- CORA-Score:} Uma métrica de maturidade científica com
propriedades matemáticas demonstráveis: monotonicidade (melhorias nunca
reduzem o score), invariância sob reescalonamento dos pesos, e
reprodutibilidade (cálculo manual confere com tracker, $\Delta = 0,003$).

\textbf{C3 --- Verificação simbólica multi-critério:} A arquitetura V1--V7
demonstra que múltiplos verificadores independentes (dimensional, algébrico,
contraexemplos, estatístico, numérico, EDO/EDP, código) produzem avaliação
mais robusta que um único critério. A correção de 44\% das falhas de modelos
Ollama via verificadores \textit{a posteriori} quantifica este ganho.

\textbf{C4 --- Geometria cognitiva:} A extensão do CORA-Eval para o espaço
dinâmico 38D (matriz de dependência, grafo cognitivo, tensor de evolução,
métrica de Fisher) introduz uma nova classe de ferramentas analíticas para
avaliação de IA --- a transição de métricas estáticas para geometria dinâmica
do aprendizado.

\textbf{C5 --- Reprodutibilidade total:} O princípio de que toda afirmação
em uma dissertação de IA deve ser reproduzível por terceiros com acesso ao
repositório, operacionalizado através de 15 itens de checklist de auditoria
(Anexo D).

\subsection{Lições Aprendidas e Recomendações para Pesquisadores}

Com base na experiência de 14 rounds de evolução e validação exaustiva,
oferecemos as seguintes recomendações para pesquisadores que desejem
implementar benchmarks científicos multidimensionais:

\textbf{R1 --- Especificação antes da implementação (SDD).} O Round 8 demonstrou
que 7 especificações modulares reduziram o tempo de correção de 3--5 iterações
para 1. Investir 20\% do tempo em especificação reduz 80\% do retrabalho.

\textbf{R2 --- Testes antes dos scores (TDD).} Nenhum score deveria ser
registrado sem um teste automatizado que o verifique. As 13 suítes TDD do
CORA-Eval (113/114 testes) são a evidência de que esta prática é viável e
produz confiança quantificável.

\textbf{R3 --- Validação externa sempre que possível.} Os 34 problemas do
Project Euler e Rosalind fornecem ground truth independente e imutável. Sempre
que possível, ancore scores em fontes externas verificáveis --- não apenas em
testes internos.

\textbf{R4 --- Documentação da reprodutibilidade.} O checklist de 15 itens
(Anexo D) deve ser preenchível por um terceiro sem conhecimento prévio do
sistema. Se um item não pode ser verificado independentemente, o score
correspondente deve ser qualificado com uma nota de incerteza.

\textbf{R5 --- Transparência sobre limitações.} A Seção 12.3 documenta 4
contraprovas --- coisas que o sistema \textit{não} consegue fazer. Esta
transparência é essencial para a credibilidade científica e deve ser prática
padrão em avaliações de IA.

\subsection{Glossário de Termos Técnicos}

\begin{table}[H]\centering\footnotesize
\caption{Glossário de siglas e termos técnicos}
\label{tab:glossario}
\begin{tabular}{p{0.18\textwidth}p{0.55\textwidth}}
\toprule
\textbf{Termo/Sigla} & \textbf{Definição} \\
\midrule
ABNT & Associação Brasileira de Normas Técnicas \\
ADR & Architecture Decision Record --- registro de decisão arquitetural \\
CORA-Eval & Cognitive ORchestrated Argumentation Evaluation \\
CORA-Score & Métrica de maturidade científica (0--4) \\
CORA-V-Score & CORA-Score ponderado por verificadores ativos \\
DCA & Dinâmica Clássica Avançada (disciplina de pós-graduação) \\
GAT & Geometric Arbitrage Theory (Farinelli, 2021) \\
MCP & Model Context Protocol --- protocolo de integração de ferramentas \\
MIRT & Multidimensional Item Response Theory \\
N1--N4 & Níveis de complexidade: Básico, Graduação, Pós-Graduação, Pesquisa \\
NFLVR & No Free Lunch with Vanishing Risk (condição de não-arbitragem) \\
SDD+TDD & Spec-Driven Development + Test-Driven Development \\
UCB1 & Upper Confidence Bound (algoritmo de seleção) \\
V1--V7 & Verificadores simbólicos do Cora-Debate \\
\bottomrule
\end{tabular}
\end{table}
